\section{A certified decay estimate for the canonical determinant}\label{sec:canonical}

\subsection{What the estimate supplies}

Let $\Delta_N=\det M_N(\theta)$, where $\theta=G+\pi^2$ and the
canonical completion is \eqref{eq:completion-basis}. The analytic
investigations \rec{R07}, \rec{R14}, \rec{R18}--\rec{R20} and
\rec{R23} culminate in the following bound.

\begin{theorem}\label{thm:canonical}
For $n=48N$,
\[
 \limsup_{N\to\infty}\frac{\log|\Delta_N|}{n^2}
 \le-1.87063289244902359295.
\]
This decay constant exceeds $A_{\rm arith}$ in
Theorem~\ref{thm:arithmetic-cost} by more than $0.00001888937116500327$.
\end{theorem}

The canonical completion retains the endpoint zero. It does not impose
the integer congruences of Theorem~\ref{thm:arithmetic-cost}. The theorem
therefore closes one analytic gap: it gives enough decay before those
additional corrections are made. Section~\ref{subsec:remaining-estimate}
states the estimate still needed to complete the contradiction.

The small margin makes the treatment of volume important. Separate
row norms lose too much information about near dependence. We retain
Gram determinants, first for the main rows and then for their endpoint
repair. A Gram determinant measures the squared volume spanned by the
vectors, and thus records this dependence.

The proof has three quantitative parts. Gram Cauchy--Schwarz bounds
$|\Delta_N|$ by the square root of the product of the row and column
Gram determinants. The row estimate splits into the main rows and the
repair required by the endpoint zero. Thus the bound for
$n^{-2}\log|\Delta_N|$ is half the sum of the first three contributions
below; changing its sign gives the decay constant.
\begin{center}
\begin{tabular}{lr}
\toprule
Main row logarithmic rate&$-0.19795874093499631465\ldots$\\
Endpoint repair, rounded upper bound&$0.00498$\\
Column logarithmic rate&$-3.54828704396305087125\ldots$\\
\midrule
Canonical decay, lower bound&$1.87063289244902359295\ldots$\\
Arithmetic cost, upper bound&$1.87061400307785858968\ldots$\\
\bottomrule
\end{tabular}
\end{center}
The resulting margin is about $1.89\times10^{-5}$ on the $n^2$ scale.
The next four subsections derive the volume decomposition and its three
bounds. The final subsection measures how much of this margin the later
integer corrections may consume.

\subsection{Separate the row and column volumes}

Put
\[
 d\nu(t)=\frac{\dd t}{\pi\sqrt{1-t^2}}\quad(-1<t<1),\qquad
 d\sigma(t)=\frac{\sqrt{1-t^2}}\pi\dd t\quad(0<t<1),
\]
and write $w(t)=t^b(1-t)^h$, $R_r=P_r/w$, $S_r=D_r/w$,
$V_r^{\rm comp}=E_r/w$. For a column polynomial $Q$, set
$\mathcal KQ(t)=\int_0^1Q(s)/(1-ts)\dd s$. The actual mixed entry is
the inner product $\pi\langle\mathcal A_r,Y_j\rangle_\nu$, where
\begin{align*}
 \mathcal A_r(t)&=|t|R_r(t)+\mathbf1_{t>0}\sqrt{1-t^2}
                 \bigl(-\tfrac32S_r(t)+24V_r^{\rm comp}(t)\bigr),\\
 Y_j(t)&=w(t)\mathcal KQ_j(t).
\end{align*}
Gram Cauchy--Schwarz consequently gives
\begin{equation}\label{eq:gram-bound}
 |\Delta_N|\le\pi^n\sqrt{\det G_{\mathcal A}\det G_Y}.
\end{equation}

The polynomial-part operation has a useful norm bound. Chebyshev
orthogonality on the full semicircle gives
$\|\Omega(R)\|_\sigma\le\|R\|_\nu$; restricting to the positive
half can only decrease the norm. Multiplication by $(1-t)^h$ does not
commute with $\Omega$, however. If $a_s(R)=\int x^sR(x)\dd\nu(x)$,
the exact defect is
\begin{equation}\label{eq:commutator}
 \Omega((1-t)^hR)=(1-t)^h\Omega(R)+J_h(R),\qquad
 J_h(R)=\sum_{\ell=1}^h(-1)^\ell\binom h\ell
                         \sum_{s=1}^\ell a_s(R)t^{\ell-s}.
\end{equation}
For $h=2$, the defect is $(a_2-2a_1)+a_1t$. Thus its dimension is
only $h$, even though the original row degrees grow much faster.

Let $\ell=56N-1$, $M=52N$, and define
\begin{align*}
 R_r(t)&=t^\ell T_{|r-g|}(1/t),\\
 V_r(t)&=\sgn(r-g)t^{\ell-1}U_{|r-g|-1}(1/t),\qquad V_g=0,\\
 C_k(t)&=\sum_{j=k}^{M-1}\binom{h+j-k-1}{j-k}t^j,\qquad 0\le k<h,\\
 a_k(x)&=\sum_{a=k+1}^h(-1)^a\binom ha x^{a-k}.
\end{align*}
Let $H_N$ be the Gram matrix of $(R_r,V_r)$ in
$L^2(\nu)\oplus L^2(\sigma)$, let $J_{rk}=\langle R_r,a_k\rangle_\nu$,
and let $B^{\rm rep}$ be the Gram matrix of the $C_k$ in $L^2(\sigma)$.
The kernel of the repair is $\sum_k a_k(x)C_k(t)$. Formula
\eqref{eq:commutator} makes its role explicit: since
$J_{rk}=[t^k]J_h(R_r)$ and
$C_k=[t^k/(1-t)^h]_{<M}$, take
\[
 \widehat E_r=w\left(\Omega(R_r)+\sum_{k=0}^{h-1}J_{rk}C_k\right).
\]
This completion agrees with $\Omega(P_r)$ in degrees $b,\ldots,L-1$,
is divisible by $w$, and has degree at most $H-1$.
Proposition~\ref{prop:completion-invariance} therefore allows us to
replace $E_r$ by $\widehat E_r$ without changing $\Delta_N$.
Since $S_r=tV_r$, the resulting row function is
\begin{align*}
 \widehat{\mathcal A}_r(t)
 &=|t|R_r(t)+24\mathbf1_{t>0}\sqrt{1-t^2}\,\Omega(R_r)(t)\\
 &\quad-\tfrac32\mathbf1_{t>0}\sqrt{1-t^2}\,tV_r(t)\\
 &\quad+24\mathbf1_{t>0}\sqrt{1-t^2}
                    \sum_{k=0}^{h-1}J_{rk}C_k(t).
\end{align*}
Use these rows in \eqref{eq:gram-bound}, with
$G_{\widehat{\mathcal A}}$ in place of $G_{\mathcal A}$.
The first line defines the bulk map from $R_r$; its norm is at most
$1+24=25$ by the contraction bound above. The maps in the next two
lines have norms at most $3/2$ and $24$, respectively.
The input triples $(R_r,V_r,\sum_kJ_{rk}C_k)$, in
$L^2(\nu)\oplus L^2(\sigma)\oplus L^2(\sigma)$, have Gram matrix
\[
 K_N=H_N+JB^{\rm rep}J^{\mathsf T},
\]
so the squared norm of the combined map is at most
$25^2+(3/2)^2+24^2=4813/4$. Consequently,
\begin{equation}\label{eq:row-repair-reduction}
 \det G_{\widehat{\mathcal A}}\le(4813/4)^n\det K_N,
 \qquad J^{\mathsf T}H_N^{-1}J\le
             (\langle a_i,a_j\rangle_\nu)_{i,j<h}.
\end{equation}
The second inequality is the contraction of an orthogonal projection.
The paired rows and their distinct lowest powers show that $H_N$ is
positive definite. All constant factors raised to $n$ disappear after
division of logarithms by $n^2$.

\subsection{The main rows: retain the spectrum of the row operator}

The paired rows split the main volume into an explicit Jacobi volume
and the extra volume contributed by a single operator. Put
$a=12N$, $m=44N$, and
\[
 \mathcal E=t^{\ell-g}\R[t]_{<g},\qquad
 \mathcal F=t^a\R[t]_{<m}.
\]
On the basis $R_d=t^\ell T_d(1/t)$, $0\le d<m$, define
$\mathcal L R_d=t^{\ell-1}U_{d-1}(1/t)$, with $\mathcal LR_0=0$.
Let $P_{\mathcal E}$ be orthogonal projection in $L^2(\sigma)$ and
$\mathcal T=(I-P_{\mathcal E})\mathcal L$. Subtracting the paired
rows gives $(0,-2V_d)$, $1\le d\le g$. Taking their Schur complement
therefore yields
\begin{equation}\label{eq:main-volume}
 \det H_N=4^g\det Z_N\det A_N^{\rm row}
                        \det(I+\mathcal T^*\mathcal T),
\end{equation}
where $Z_N$ and $A_N^{\rm row}$ are the Gram matrices of these $V_d$
and of $R_0,\ldots,R_{m-1}$, respectively.

The first two factors can be evaluated through monic Jacobi norms.
For weight $x^u(1-x)^v$ on $[0,1]$, that norm is
\[
 \frac{j!\,\Gamma(j+u+1)\Gamma(j+v+1)\Gamma(j+u+v+1)}
 {(2j+u+v+1)\Gamma(2j+u+v+1)^2}.
\]
The triangular leading coefficients of the Chebyshev bases contribute
$2^{(m-1)(m-2)}$ and $2^{g(g-1)}$. Replacing the remaining bounded
weight factors changes the logarithm by $O(n)$. Stirling's formula
then gives the limiting base contribution
\begin{align*}
 b_0&=\tfrac{61}{72}\log2+2\mathcal I(1/4,0,11/24)
                         +\mathcal I(13/6,0,1/12)\\
 &=-0.243039538134996314659\ldots,\\
 \mathcal I(u,v,T)&=\int_0^T
 \{f(t)+f(t+u)+f(t+v)+f(t+u+v)-2f(2t+u+v)\}\dd t,
\end{align*}
Here $f(x)=x\log x$, with $f(0)=0$.

We still need the last factor in \eqref{eq:main-volume}. If
$s_1\ge\cdots\ge s_m\ge0$ are the singular values of $\mathcal T$, then
\[
 \log\det(I+\mathcal T^*\mathcal T)
                =\sum_{j=1}^m\log(1+s_j^2).
\]
Bounding every $s_j$ by $s_1$ would charge the worst direction $m$ times.
Instead, we approximate $\mathcal T$ by operators of successively larger
rank. The error after a rank-$l$ approximation bounds $s_{l+1}$, allowing
the later, smaller directions to contribute less to this sum.

For $R\in\mathcal F$,
put $q_R(x)=x^\ell R(1/x)$. The useful exact representation is
\[
 \mathcal LR(t)=\left[\frac{R(t)}{\sqrt{1-t^2}}\right]_{<\ell}
 =\frac{R(t)}{\sqrt{1-t^2}}
   -t^\ell\int\frac{q_R(x)}{1-xt}\dd\nu(x).
\]
The first term is a contraction from $L^2(\nu)$ to $L^2(\sigma)$.
Approximate the second term by rational functions. Choose $g+l$ zeros
$z_i\in(0,1)$ and $l$ poles $y_j\in(-1,1)$, and set
\[
 P(t)=\prod_i(t-z_i),\quad Q(t)=\prod_j(1-y_jt),\quad
 W_*(t)=t^{\ell-g}P(t)/Q(t),\quad
 B_*(x)=\frac{\prod_j(x-y_j)}{\prod_i(1-z_ix)}.
\]
Polynomial division gives
\[
 \frac{t^\ell}{1-xt}
 =\frac{t^{\ell-g}v_x(t)}{Q(t)}
       +\frac{W_*(t)B_*(x)}{1-xt},\qquad \deg v_x<g+l.
\]
The first space contains $\mathcal E$, so it has rank at most $l$
after projection. This converts a uniform bound on the second term
into a bound on the $(l+1)$st singular value.

The necessary evaluation estimate is
\[
 |q_R(x)|\le C_N e^{n\mathcal E_0(x)}\|R\|_\nu,
 \qquad C_N=O(n),
\]
where, with $\tau=11/24$ and $\alpha=1/4$,
\[
 \mathcal E_0(x)=\inf_{u>0}
 \{\tau\log(1+u)+(\tau+\alpha)\log(1+(1-x^2)u)-\tau\log u\}.
\]
At $x=\pm1$ its value is zero. To obtain the estimate, expand in the
orthonormal Jacobi basis and use the positive coefficient sum
$\sum_i\binom K i\binom{K+a}i(1-x^2)^i$, $K=22N$.
It is the coefficient of $u^K$ in
$(1+u)^K(1+(1-x^2)u)^{K+a}$; bounding that coefficient at positive
$u$ and minimizing gives the displayed field. The evaluation norms
supply, for example, $C_N=\sqrt{\pi m(2K+a+2)}$.

For $l/n\to r$, describe the zeros and poles by positive measures
$p,q$ of masses $1/12+r,r$. The two potentials to bound are
\begin{align*}
 \mathcal U(t)&=\tfrac{13}{12}\log t
       +\int\log|t-z|\dd p(z)-\int\log(1-yt)\dd q(y),\\
 \mathcal V(x)&=\mathcal E_0(x)
       +\int\log|x-y|\dd q(y)-\int\log(1-zx)\dd p(z).
\end{align*}
If their upper bounds sum to $e(r)\ge0$, the singular-value estimate
is $s_{l+1}(\mathcal T)\le1+\exp(n(e(r)+o(1)))$.
The accompanying rational atomic measures certify the following values.
\begin{center}\small
\begin{tabular}{rrrrrrr}
\toprule
$r$&0&.02&.04&.06&.08&.10\\
$e(r)$&.43055090&.32535779&.24085716&.16822647&.10576227&.05324220\\
\midrule
$r$&.12&.14&.16&.18&.20&.22\\
$e(r)$&.01233031&.00066131&.00029828&.00023083&.00019427&.00015768\\
\midrule
$r$&.24&&&&&\\
$e(r)$&.00012891&&&&&\\
\bottomrule
\end{tabular}
\end{center}
Convex interpolation of the measures proves the linear interpolant is
an upper bound between nodes. After $.24$, adding equal masses at
$1/2$ to both measures cannot increase either potential, since
$|(x-1/2)/(1-x/2)|\le1$ on $[-1,1]$. This extends the final bound
to $r=11/12$. Rational masses can be rounded to integer multiplicities
with only $O(1)$ residual nodes. Upper bounds for the logarithmic
potentials persist in this passage; zeros make the potential smaller,
not larger. At an index $j$ with $j/n\approx r$, the contribution
$\log(1+s_j^2)$ is bounded on the leading scale by $2ne(r)$.
Summing these contributions over the $m=11n/12$ singular values gives
the integral $2\int_0^{11/12}e(r)\dd r$ after division by $n^2$.
Together with the Jacobi contribution in \eqref{eq:main-volume}, this gives
\begin{equation}\label{eq:main-row-bound}
 \limsup\frac{\log\det H_N}{n^2}
 \le b_0+2\int_0^{11/12}e(r)\dd r
 =b_0+\frac{112701993}{2500000000}
 =-0.19795874093499631465\ldots.
\end{equation}
The certificate uses $192$-bit interval arithmetic and a global
curvature bound on each subinterval; it checks both potentials at all
13 nodes, rather than sampling their maxima.

\subsection{The repair: keep the negative interaction between degrees}

The repair in \eqref{eq:commutator} is already part of the canonical
matrix. It is distinct from the later arithmetic correction $UW$.
The kernel simplifies to
\[
 K(x,t)=-x\sum_{j=0}^{h-1}(1-x)^j
                       \sum_{r=0}^{M-1}\binom{r+j}{j}t^r.
\]
Let $A_{ij}=\langle-x(1-x)^i,-x(1-x)^j\rangle_\nu$ and
$B_{ij}=\sum_{r=0}^{M-1}\binom{r+i}i\binom{r+j}j$.
The elementary bound $\|\sum c_rt^r\|_\sigma^2\le M\sum c_r^2$
and \eqref{eq:row-repair-reduction} give
\[
 \det K_N/\det H_N\le\det(I_h+MAB).
\]
For a $k$-element set $I$, let $A_I$ denote its principal minor.
Cauchy--Binet and Gram Cauchy--Schwarz imply
\begin{equation}\label{eq:compound-bound}
 e_k(AB)\le\binom hk^2\max_{|I|=k}\det A_I\det B_I.
\end{equation}
Here $e_k(AB)$ is the $k$th elementary symmetric function of the
eigenvalues of $AB$, so $\det(I_h+MAB)=\sum_{k=0}^h M^k e_k(AB)$.
This keeps the interaction between the selected degrees. Replacing
the minors by separate vector norms would discard that saving.

Put $L_*=M+h=27h$ and $\mathcal H_I=(1/(i+j+1))_{i,j\in I}$.
The two comparisons needed are
\begin{equation}\label{eq:repair-minors}
 \det A_I\le h^{2k}2^{2\sum_I j}\det\mathcal H_I,\qquad
 \det B_I\le(Mh^2)^k
       \frac{L_*^{2\sum_I j}}{\prod_I j!^2}\det\mathcal H_I.
\end{equation}
For the first, change variables $y=1-x$ and use $x^2\le1$.
The polynomial evaluation inequality
$\sup_{[0,L_0]}|f|^2\le(h^2/L_0)\int_0^{L_0}|f|^2$
for $\deg f<h$, obtained from shifted Legendre polynomials, bounds
the arcsine Gram by the uniform Gram on $[0,2]$.

For the second comparison, the factorial Schur tableau identity
\cite[Theorem~2 and (7)--(8), (14)]{BMN} gives, for increasing
nonnegative $r_a$ and increasing $j_b<h$,
\[
 0\le\det\left[\prod_{v=1}^{j_b}(r_a+v)\right]
       \le\det[(r_a+h)^{j_b}].
\]
Indeed, divide by the Vandermonde and use the partition
$\lambda_i=j_{k+1-i}-k+i$. Every tableau factor has shift between
$1$ and $j_k<h$, so replacing it by $h$ increases every nonnegative
summand. Multiplication by the Vandermonde proves the comparison.
Apply it to each evaluation minor of $B_I$, divide column $j$ by
$j!$, and then use the polynomial evaluation inequality on $[0,L_*]$.
This proves \eqref{eq:repair-minors}. The binomial comparison is a
comparison of determinants; no entrywise matrix ordering is asserted.

The saving from these minors is visible with just two degrees. For
$I=\{0,1\}$,
\[
 \mathcal H_I=\begin{pmatrix}1&1/2\\1/2&1/3\end{pmatrix},\qquad
 \det\mathcal H_I=\frac1{12},\qquad
 (\mathcal H_I)_{00}(\mathcal H_I)_{11}=\frac13.
\]
The determinant is four times smaller than the product of the separate
squared norms. For general $i<j$, the corresponding ratio is
\[
 \frac{\det\mathcal H_{\{i,j\}}}
      {(\mathcal H_{\{i,j\}})_{ii}(\mathcal H_{\{i,j\}})_{jj}}
       =\left(\frac{j-i}{i+j+1}\right)^2.
\]
Nearby powers are nearly dependent, and their joint volume records this.
Both comparisons in \eqref{eq:repair-minors} contain the same Cauchy
minor, so their product contributes the factor
$((j-i)/(i+j+1))^4$. These factors reduce the bound; their logarithms
give the negative interactions retained below.

The Cauchy determinant formula now gives
\begin{align*}
 \det K_N/\det H_N&\le(4M^2h^4)^h\max_{I\subset\{0,\ldots,h-1\}}\mathcal W_I,\\
 \mathcal W_I&=\prod_{j\in I}\left(\frac{(2L_*)^j}{j!}\right)^2
     \prod_{i<j\in I}\left|\frac{i-j}{i+j+1}\right|^4.
\end{align*}
Use the midpoint lattice $x_j=(j+1/2)/h$. Stirling's formula turns
the first product into the field
$w_0(x)=2x(\log(54/x)+1)$. The pair interaction is
$G_0(x,y)=\log|(x-y)/(x+y)|$.

Its Chebyshev expansion is
$G_0=-4\sum_{k\ge1,\ k\ \mathrm{odd}}T_k(x)T_k(y)/k$.
To make the passage to a uniform bound legitimate, damp by $\rho^k$,
$0<\rho<1$. The damped kernel is
\[
 G_\rho(x,y)=\tfrac12\log\frac{D_0-D_1}{D_0+D_1},\quad
 D_0=(1-\rho^2)^2+4\rho^2(x^2+y^2),\quad
 D_1=4\rho(1+\rho^2)xy.
\]
On $[0,1]^2$, $G_0\le G_\rho$. After cross multiplication the
required difference is
$4xy(1-\rho)^2((1+\rho)^2-2\rho(x^2+y^2))\ge0$.
To expose the squares in this interaction, define the averaged moments
\[
 m_k(I)=\frac1h\sum_{j\in I}T_k(x_j),\qquad k\ge1\text{ odd}.
\]
The field contribution is at most $h\sum_{j\in I}w_0(x_j)$, using
$j!\ge(j/e)^j$ and the monotonicity of $w_0$ on $[0,1]$.
Replacing $G_0$ by $G_\rho$ and inserting the diagonal in the pair
sum gives
\[
 \frac{\log\mathcal W_I}{h^2}
 \le\frac1h\sum_{j\in I}w_0(x_j)
       -8\sum_{k\ {\rm odd}}\frac{\rho^km_k(I)^2}{k}
       +\frac4h\log\frac{1+\rho}{1-\rho}.
\]
Here the last term removes the inserted diagonal:
$G_\rho(x,x)\ge-2\log((1+\rho)/(1-\rho))$.
It tends to zero for fixed $\rho$ as $h\to\infty$.

Choose any finitely supported real sequence $v_k$ on the positive
odd indices. The elementary inequality
\[
 -m_k(I)^2\le v_k^2-2v_km_k(I)
\]
replaces each negative square by a term linear in the moments.
After substituting their definition, site $x_j$ contributes the field
$w_0(x_j)-16\sum_{k\ {\rm odd}}\rho^kv_kT_k(x_j)/k$ if it belongs
to $I$, and zero otherwise. Its contribution is therefore at most
the positive part of this field. Summing over all sites removes the
dependence on $I$. First let $h\to\infty$, so the full lattice sum
becomes an integral, and then let $\rho\uparrow1$. Since $n=24h$,
the resulting bound is
\begin{equation}\label{eq:repair-dual}
 \limsup\frac{\log(\det K_N/\det H_N)}{n^2}
 \le\frac1{576}\left\{
 8\sum_{k\ \mathrm{odd}}\frac{v_k^2}{k}
 +\int_0^1\left(w_0(x)-16\sum_{k\ \mathrm{odd}}\frac{v_kT_k(x)}k\right)_+\dd x
 \right\}.
\end{equation}
This order of limits avoids using an undamped logarithmic series on
the diagonal. It also retains the occupancy restriction on $I$.

The supplied trial uses odd indices through $79$ and rational
coefficients with denominator $10^{12}$. A $512$-bit calculation on
$16384$ intervals bounds the positive-part integral by
$0.85979724477769159413$. On each interval the polynomial's second
derivative is bounded by
$\frac{16}{3}\sum |v_k|k(k^2-1)$; the field's second derivative is
negative. An affine upper bound, with this quadratic error added,
can be integrated exactly even where it crosses zero. The result is
\begin{equation}\label{eq:repair-rate}
 \limsup\frac{\log(\det K_N/\det H_N)}{n^2}
 \le0.004978210597652675202<0.00498.
\end{equation}

\subsection{The columns: an energy bound with certified maxima}

Expanding the column Gram determinant twice with the Cauchy determinant
formula expresses it as a positive integral:
\[
 \det G_Y=\frac1{(n!)^3}\int
 \frac{V(t)^2V(s)^2V(u)^2
       \prod_i w(t_i)^2s_i^b(1-s_i)^qu_i^b(1-u_i)^q}
      {\prod_{i,j}(1-t_is_j)(1-t_iu_j)}
 \prod_i\dd\nu(t_i)\dd s_i\dd u_i,
\]
where $V$ denotes the Vandermonde, $t_i\in(-1,1)$ and $s_i,u_i\in(0,1)$.
This positivity permits a uniform bound on the integrand.
Make the substitution $t=2x/(1+x^2)$, and put
\[
 \beta=7/48,\quad\eta=1/24,\quad\gamma=1/12,
\]
\[
 W_0(x)=\beta\log|x|+2\eta\log(1-x)-(\beta+\eta)\log(1+x^2),
 \qquad Z_0(s)=\beta\log s+\gamma\log(1-s).
\]
For any finite sequences $p_k,v_k$, the energy calculation gives
\begin{align}\label{eq:column-dual}
 \limsup\frac{\log\det G_Y}{2n^2}
 \le{}&(\beta-1)\log2+\sum_k\frac{p_k^2}{k}
                         +\frac12\sum_k\frac{v_k^2}{k}\nonumber\\
 &+\sup_{-1\le x\le1}\left(W_0(x)-2\sum_k\frac{p_kT_k(x)}k
                                      -\sum_k\frac{v_kx^k}k\right)\nonumber\\
 &+\sup_{0\le s\le1}\left(Z_0(s)+2\sum_k\frac{v_kT_k(s)}k\right).
\end{align}

For completeness, the square completion behind this estimate keeps
three empirical measures, one each for $x,s,u$. Let $A_k$ be the
Chebyshev moments of the first, $B_k$ its power moments, and $C_k,D_k$
the Chebyshev moments of the other two. The interaction, apart from
its constant, is
\[
 -\sum_k\frac{A_k^2}{k}
 -\frac14\sum_k\frac{(B_k-2C_k)^2+(B_k-2D_k)^2}{k}.
\]
Applying $-a^2\le p^2-2pa$ gives \eqref{eq:column-dual}.
To justify the expansion uniformly at endpoints and repeated points,
use dampings $\tau=1-\varepsilon$ for the $x$ Chebyshev terms,
$\rho(x)=1-\varepsilon(1-x)$ for its power terms, and
$\sigma(s)=1-\varepsilon\sqrt{1-s}$ for the $s,u$ terms.
The kernel comparison costs $O(\varepsilon)$ on the $n^2$ scale;
diagonal errors are $O_\varepsilon(n)$, and the endpoint logarithms
are absorbed by the positive exponents $\beta,\eta,\gamma$.
Let $n\to\infty$ first, then $\varepsilon\downarrow0$. This yields
the displayed bound for the full integral, not just its interior.

For the rational degree-$200$ trial, the two maxima are bounded by
\[
 -\frac{1798431303}{5000000000},\qquad
 -\frac{7455295417}{5000000000},
\]
respectively. Clearing the denominators in the derivatives gives exact
polynomials. Sturm counts give $96$ roots on $(-1,0)$, $89$ on $(0,1)$
for the first function, and $103$ on $(0,1)$ for the second.
The certificate encloses all $288$ roots, evaluates every critical
value and the regular endpoint $x=-1$, and handles the other endpoints
by their logarithmic limits. It proves
\begin{equation}\label{eq:column-rate}
 \limsup\frac{\log\det G_Y}{n^2}
 \le-3.54828704396305087125.
\end{equation}

\subsection{The margin and the estimate still missing}\label{subsec:remaining-estimate}

Combining \eqref{eq:gram-bound}, \eqref{eq:row-repair-reduction},
\eqref{eq:main-row-bound}, \eqref{eq:repair-rate} and
\eqref{eq:column-rate} proves Theorem~\ref{thm:canonical}.
Taking half the sum of the three contributions listed at the start of
the section gives the stated decay bound.
The unrounded repair bound gives a decay constant of
$1.87063378715019725535\ldots$. Its margin over the arithmetic cost is
$0.00001978407233866567\ldots$.

Where $M_N(\theta)$ is invertible, the determinant lemma writes the
effect of a later integer correction as
\[
 P_{N,U}(\theta)=\Delta_N
              \det(I_b+WM_N(\theta)^{-1}U).
\]
For example, with the induced maximum-row-sum norm, it would suffice
to choose admissible corrections along the prime scales such that
\begin{equation}\label{eq:missing-correction-bound}
 \limsup\frac1n\log\bigl(1+\|WM_N(\theta)^{-1}U\|_\infty\bigr)
 <0.00012952711656002246.
\end{equation}
Indeed, the logarithm of the extra determinant is at most $b$ times
this logarithm, and $b/n=7/48$. The unrounded repair allows the slightly
larger threshold $0.00013566221032227892$. These are sufficient
conditions, not estimates that have been proved. One could also bound
the corrected determinant directly and avoid the inverse.

Appendix~\ref{app:audit} records the independent audit's scope.
The next section examines ways to obtain the missing correction estimate
and the limitations established for those approaches.
